\chapter{Fonctions : définir, représenter, modéliser}
\section*{Exercice 12.1 -- Images et antécédents}
$f(-2)=9$. Les antécédents de $4$ vérifient $x-1=\pm2$, donc sont $-1$ et $3$. $0$ possède le seul antécédent $1$ ; $-1$ n'en possède aucun puisque la fonction est non négative.

\section*{Exercice 12.2 -- Domaine}
La racine exige $x\leq3$, le dénominateur $x\ne-2$. Le domaine est $]-\infty,-2[\cup]-2,3]$.

\section*{Exercice 12.3 -- Appartenance}
Pour $A$, $2^2+1=5$, donc oui. Pour $B$, $(-2)^2+1=5\ne-3$, donc non. Il faut remplacer l'abscisse dans la formule et comparer à l'ordonnée.

\section*{Exercice 12.4 -- Tableau insuffisant}
Non : $f(x)=0$ et $g(x)=x(x-1)(x+1)$ coïncident sur ces trois entrées, mais $g(2)=6\ne0=f(2)$. Un tableau fini ne détermine pas une fonction quelconque sur $\R$.

\section*{Exercice 12.5 -- Modèle géométrique}
Les côtés vérifient $2x+2y=30$, donc $y=15-x$. L'aire est $A(x)=x(15-x)$ pour $0<x<15$. Le domaine du modèle est plus restreint que celui du polynôme.

\section*{Exercice 12.6 -- Modèle discret}
$u(4)=18$. L'équation $8+2{,}5n=19$ donne $n=4{,}4$, qui n'est pas entier. Ce prix n'est donc pas réalisable dans le modèle, même s'il existe un point d'abscisse $4{,}4$ sur la droite prolongée.
